% Generated by analysis/apa_refs.py from refs.bib (APA 7). Do not edit by hand.
\begin{thebibliography}{48}
\bibitem[{Aures(1985)Aures}]{aures1985}
Aures, W. (1985). Berechnungsverfahren für den sensorischen Wohlklang beliebiger Schallsignale. \emph{Acustica}, \emph{59}(2), 130--141. \url{https://www.ingentaconnect.com/content/dav/aaua/1985/00000059/00000002/art00008}

\bibitem[{Bowling \apaand{} Purves(2015)Bowling and Purves}]{bowling2015}
Bowling, D. L., \& Purves, D. (2015). A biological rationale for musical consonance. \emph{Proceedings of the National Academy of Sciences}, \emph{112}(36), 11155--11160. \url{https://doi.org/10.1073/pnas.1505768112}

\bibitem[{Bowling et~al.(2018)Bowling, Purves, and Gill}]{bowling2018}
Bowling, D. L., Purves, D., \& Gill, K. Z. (2018). Vocal similarity predicts the relative attraction of musical chords. \emph{Proceedings of the National Academy of Sciences}, \emph{115}(1), 216--221. \url{https://doi.org/10.1073/pnas.1713206115}

\bibitem[{Brown(1910)Brown}]{brown1910}
Brown, W. (1910). Some experimental results in the correlation of mental abilities. \emph{British Journal of Psychology}, \emph{3}(3), 296--322. \url{https://doi.org/10.1111/j.2044-8295.1910.tb00207.x}

\bibitem[{Burgoyne et~al.(2011)Burgoyne, Wild, and Fujinaga}]{burgoyne2011}
Burgoyne, J. A., Wild, J., \& Fujinaga, I. (2011). An expert ground truth set for audio chord recognition and music analysis. In \emph{Proceedings of the 12th International Society for Music Information Retrieval Conference} (pp. 633--638). \url{https://ismir2011.ismir.net/papers/OS8-1.pdf}

\bibitem[{De Roure(2026)De Roure}]{deroure2026}
De Roure, D. (2026). \emph{An asymmetric formula for interval consonance and its relation to harmonic coincidence} [Preprint]. arXiv. \url{https://doi.org/10.48550/arXiv.2606.16412}

\bibitem[{DerSimonian \apaand{} Laird(1986)DerSimonian and Laird}]{dersimonian1986}
DerSimonian, R., \& Laird, N. (1986). Meta-analysis in clinical trials. \emph{Controlled Clinical Trials}, \emph{7}(3), 177--188. \url{https://doi.org/10.1016/0197-2456(86)90046-2}

\bibitem[{Digital \apaand{} Cognitive Musicology Lab, EPFL(2025)Digital and Cognitive Musicology Lab, EPFL}]{dcmlab}
Digital and Cognitive Musicology Lab, EPFL. (2025). \emph{consonance: Collected consonance rating datasets} [Data set]. \url{https://github.com/DCMLab/consonance}

\bibitem[{Eerola \apaand{} Lahdelma(2021)Eerola and Lahdelma}]{eerola2021}
Eerola, T., \& Lahdelma, I. (2021). The anatomy of consonance/dissonance: Evaluating acoustic and cultural predictors across multiple datasets with chords. \emph{Music \& Science}, \emph{4}. \url{https://doi.org/10.1177/20592043211030471}

\bibitem[{Eerola \apaand{} Lahdelma(2022)Eerola and Lahdelma}]{eerola2022register}
Eerola, T., \& Lahdelma, I. (2022). Register impacts perceptual consonance through roughness and sharpness. \emph{Psychonomic Bulletin \& Review}, \emph{29}(3), 800--808. \url{https://doi.org/10.3758/s13423-021-02033-5}

\bibitem[{Fastl \apaand{} Zwicker(2007)Fastl and Zwicker}]{zwicker2007}
Fastl, H., \& Zwicker, E. (2007). \emph{Psychoacoustics: Facts and models} (3rd ed.). Springer. \url{https://doi.org/10.1007/978-3-540-68888-4}

\bibitem[{Harrison(2019)Harrison}]{incondata}
Harrison, P. M. C. (2019). \emph{inconData: Consonance perception datasets} (Version 0.1.0) [Data set]. \url{https://github.com/pmcharrison/inconData}

\bibitem[{Harrison(2023)Harrison}]{marjieh2024code}
Harrison, P. M. C. (2023). \emph{Data and analysis code for Marjieh et al. (2024)} [Data set and computer code]. \url{https://gitlab.com/pmcharrison/timbre-and-consonance-paper}

\bibitem[{Harrison(2024)Harrison}]{incon}
Harrison, P. M. C. (2024). \emph{incon: Computational models of simultaneous consonance} (Version 0.5.0) [Computer software]. \url{https://github.com/pmcharrison/incon}

\bibitem[{Harrison \apaand{} MacConnachie(2024)Harrison and MacConnachie}]{harrison2024carillon}
Harrison, P. M. C., \& MacConnachie, J. (2024). Consonance in the carillon. \emph{The Journal of the Acoustical Society of America}, \emph{156}(2), 1111--1122. \url{https://doi.org/10.1121/10.0028167}

\bibitem[{Harrison \apaand{} Pearce(2018)Harrison and Pearce}]{harrison2018}
Harrison, P. M. C., \& Pearce, M. T. (2018). \emph{Dissociating sensory and cognitive theories of harmony perception through computational modeling} [Preprint]. PsyArXiv. \url{https://doi.org/10.31234/osf.io/wgjyv}

\bibitem[{Harrison \apaand{} Pearce(2020)Harrison and Pearce}]{harrison2020}
Harrison, P. M. C., \& Pearce, M. T. (2020). Simultaneous consonance in music perception and composition. \emph{Psychological Review}, \emph{127}(2), 216--244. \url{https://doi.org/10.1037/rev0000169}

\bibitem[{Hartung \apaand{} Knapp(2001)Hartung and Knapp}]{hartung2001}
Hartung, J., \& Knapp, G. (2001). A refined method for the meta-analysis of controlled clinical trials with binary outcome. \emph{Statistics in Medicine}, \emph{20}(24), 3875--3889. \url{https://doi.org/10.1002/sim.1009}

\bibitem[{Higgins \apaand{} Thompson(2002)Higgins and Thompson}]{higgins2002}
Higgins, J. P. T., \& Thompson, S. G. (2002). Quantifying heterogeneity in a meta-analysis. \emph{Statistics in Medicine}, \emph{21}(11), 1539--1558. \url{https://doi.org/10.1002/sim.1186}

\bibitem[{Holm(1979)Holm}]{holm1979}
Holm, S. (1979). A simple sequentially rejective multiple test procedure. \emph{Scandinavian Journal of Statistics}, \emph{6}(2), 65--70. \url{https://www.jstor.org/stable/4615733}

\bibitem[{Hutchinson \apaand{} Knopoff(1978)Hutchinson and Knopoff}]{hutchinson1978}
Hutchinson, W., \& Knopoff, L. (1978). The acoustic component of Western consonance. \emph{Interface}, \emph{7}(1), 1--29. \url{https://doi.org/10.1080/09298217808570246}

\bibitem[{IntHout et~al.(2014)IntHout, Ioannidis, and Borm}]{inthout2014}
IntHout, J., Ioannidis, J. P. A., \& Borm, G. F. (2014). The Hartung-Knapp-Sidik-Jonkman method for random effects meta-analysis is straightforward and considerably outperforms the standard DerSimonian-Laird method. \emph{BMC Medical Research Methodology}, \emph{14}, 25. \url{https://doi.org/10.1186/1471-2288-14-25}

\bibitem[{Johnson-Laird et~al.(2012)Johnson-Laird, Kang, and Leong}]{johnsonlaird2012}
Johnson-Laird, P. N., Kang, O. E., \& Leong, Y. C. (2012). On musical dissonance. \emph{Music Perception}, \emph{30}(1), 19--35. \url{https://doi.org/10.1525/mp.2012.30.1.19}

\bibitem[{Kaklamani \apaand{} Simserides(2026)Kaklamani and Simserides}]{kaklamani2025}
Kaklamani, S., \& Simserides, C. (2026). Psychoacoustic study of simple-tone dyads: Frequency ratio and pitch. \emph{Acoustics}, \emph{8}(1), 14. \url{https://doi.org/10.3390/acoustics8010014}

\bibitem[{Kameoka \apaand{} Kuriyagawa(1969)Kameoka and Kuriyagawa}]{kameoka1969}
Kameoka, A., \& Kuriyagawa, M. (1969). Consonance theory part II: Consonance of complex tones and its calculation method. \emph{The Journal of the Acoustical Society of America}, \emph{45}(6), 1460--1469. \url{https://doi.org/10.1121/1.1911624}

\bibitem[{Lahdelma et~al.(2022)Lahdelma, Eerola, and Armitage}]{lahdelma2022}
Lahdelma, I., Eerola, T., \& Armitage, J. (2022). Is harmonicity a misnomer for cultural familiarity in consonance preferences? \emph{Frontiers in Psychology}, \emph{13}, 802385. \url{https://doi.org/10.3389/fpsyg.2022.802385}

\bibitem[{Marjieh et~al.(2024)Marjieh, Harrison, Lee, Deligiannaki, and Jacoby}]{marjieh2024}
Marjieh, R., Harrison, P. M. C., Lee, H., Deligiannaki, F., \& Jacoby, N. (2024). Timbral effects on consonance disentangle psychoacoustic mechanisms and suggest perceptual origins for musical scales. \emph{Nature Communications}, \emph{15}, 1482. \url{https://doi.org/10.1038/s41467-024-45812-z}

\bibitem[{Masina et~al.(2022)Masina, Lo Presti, and Stanzial}]{masina2022}
Masina, I., Lo Presti, G., \& Stanzial, D. (2022). Dyad's consonance and dissonance: Combining the compactness and roughness approaches. \emph{The European Physical Journal Plus}, \emph{137}(11), 1254. \url{https://doi.org/10.1140/epjp/s13360-022-03456-2}

\bibitem[{McBride(2025)McBride}]{mcbride2025}
McBride, J. M. (2025). \emph{Musical consonance: A review of theory and evidence on perception and preference of auditory roughness in humans and other animals} [Preprint]. arXiv. \url{https://doi.org/10.48550/arXiv.2510.14159}

\bibitem[{McDermott et~al.(2010)McDermott, Lehr, and Oxenham}]{mcdermott2010}
McDermott, J. H., Lehr, A. J., \& Oxenham, A. J. (2010). Individual differences reveal the basis of consonance. \emph{Current Biology}, \emph{20}(11), 1035--1041. \url{https://doi.org/10.1016/j.cub.2010.04.019}

\bibitem[{McDermott et~al.(2016)McDermott, Schultz, Undurraga, and Godoy}]{mcdermott2016}
McDermott, J. H., Schultz, A. F., Undurraga, E. A., \& Godoy, R. A. (2016). Indifference to dissonance in native Amazonians reveals cultural variation in music perception. \emph{Nature}, \emph{535}(7613), 547--550. \url{https://doi.org/10.1038/nature18635}

\bibitem[{Milne et~al.(2023)Milne, Smit, Sarvasy, and Dean}]{milne2023}
Milne, A. J., Smit, E. A., Sarvasy, H. S., \& Dean, R. T. (2023). Evidence for a universal association of auditory roughness with musical stability. \emph{PLOS ONE}, \emph{18}(9), e0291642. \url{https://doi.org/10.1371/journal.pone.0291642}

\bibitem[{Parncutt(1988)Parncutt}]{parncutt1988}
Parncutt, R. (1988). Revision of Terhardt's psychoacoustical model of the root(s) of a musical chord. \emph{Music Perception}, \emph{6}(1), 65--93. \url{https://doi.org/10.2307/40285416}

\bibitem[{Parncutt(1989)Parncutt}]{parncutt1989}
Parncutt, R. (1989). \emph{Harmony: A psychoacoustical approach}. Springer. \url{https://doi.org/10.1007/978-3-642-74831-8}

\bibitem[{Parncutt \apaand{} Hair(2011)Parncutt and Hair}]{parncutt2011}
Parncutt, R., \& Hair, G. (2011). Consonance and dissonance in music theory and psychology: Disentangling dissonant dichotomies. \emph{Journal of Interdisciplinary Music Studies}, \emph{5}(2), 119--166. \url{https://web.archive.org/web/20210227092804/http://musicstudies.org/wp-content/uploads/2017/01/Parncutt_JIMS_11050202.pdf}

\bibitem[{Parncutt et~al.(2019)Parncutt, Reisinger, Fuchs, and Kaiser}]{parncutt2019}
Parncutt, R., Reisinger, D., Fuchs, A., \& Kaiser, F. (2019). Consonance and prevalence of sonorities in Western polyphony: Roughness, harmonicity, familiarity, evenness, diatonicity. \emph{Journal of New Music Research}, \emph{48}(1), 1--20. \url{https://doi.org/10.1080/09298215.2018.1477804}

\bibitem[{Philharmonia Orchestra(n.d.)Philharmonia Orchestra}]{philharmonia}
Philharmonia Orchestra. (n.d.). \emph{Sound samples}. Retrieved August 2026, from \url{https://philharmonia.co.uk/resources/sound-samples/}

\bibitem[{Plomp \apaand{} Levelt(1965)Plomp and Levelt}]{plomp1965}
Plomp, R., \& Levelt, W. J. M. (1965). Tonal consonance and critical bandwidth. \emph{The Journal of the Acoustical Society of America}, \emph{38}(4), 548--560. \url{https://doi.org/10.1121/1.1909741}

\bibitem[{Popescu et~al.(2019)Popescu, Neuser, Neuwirth, Bravo, Mende, Boneh, Moss, and Rohrmeier}]{popescu2019}
Popescu, T., Neuser, M. P., Neuwirth, M., Bravo, F., Mende, W., Boneh, O., Moss, F. C., \& Rohrmeier, M. (2019). The pleasantness of sensory dissonance is mediated by musical style and expertise. \emph{Scientific Reports}, \emph{9}, 1070. \url{https://doi.org/10.1038/s41598-018-35873-8}

\bibitem[{Pressnitzer et~al.(2001)Pressnitzer, Patterson, and Krumbholz}]{pressnitzer2001}
Pressnitzer, D., Patterson, R. D., \& Krumbholz, K. (2001). The lower limit of melodic pitch. \emph{The Journal of the Acoustical Society of America}, \emph{109}(5), 2074--2084. \url{https://doi.org/10.1121/1.1359797}

\bibitem[{Sethares(1993)Sethares}]{sethares1993}
Sethares, W. A. (1993). Local consonance and the relationship between timbre and scale. \emph{The Journal of the Acoustical Society of America}, \emph{94}(3), 1218--1228. \url{https://doi.org/10.1121/1.408175}

\bibitem[{Sethares(2005)Sethares}]{sethares2005}
Sethares, W. A. (2005). \emph{Tuning, timbre, spectrum, scale} (2nd ed.). Springer. \url{https://doi.org/10.1007/b138848}

\bibitem[{Shamma \apaand{} Klein(2000)Shamma and Klein}]{shamma2000}
Shamma, S., \& Klein, D. (2000). The case of the missing pitch templates: How harmonic templates emerge in the early auditory system. \emph{The Journal of the Acoustical Society of America}, \emph{107}(5), 2631--2644. \url{https://doi.org/10.1121/1.428649}

\bibitem[{Spearman(1910)Spearman}]{spearman1910}
Spearman, C. (1910). Correlation calculated from faulty data. \emph{British Journal of Psychology}, \emph{3}(3), 271--295. \url{https://doi.org/10.1111/j.2044-8295.1910.tb00206.x}

\bibitem[{Stolzenburg(2015)Stolzenburg}]{stolzenburg2015}
Stolzenburg, F. (2015). Harmony perception by periodicity detection. \emph{Journal of Mathematics and Music}, \emph{9}(3), 215--238. \url{https://doi.org/10.1080/17459737.2015.1033024}

\bibitem[{Terhardt(1974)Terhardt}]{terhardt1974}
Terhardt, E. (1974). Pitch, consonance, and harmony. \emph{The Journal of the Acoustical Society of America}, \emph{55}(5), 1061--1069. \url{https://doi.org/10.1121/1.1914648}

\bibitem[{Tonalsoft(n.d.)Tonalsoft}]{erlich}
Tonalsoft. (n.d.). \emph{Harmonic entropy}. Retrieved September 2026, from \url{http://www.tonalsoft.com/enc/h/harmonic-entropy.aspx}

\bibitem[{University of Iowa Electronic Music Studios(n.d.)University of Iowa Electronic Music Studios}]{iowamis}
University of Iowa Electronic Music Studios. (n.d.). \emph{Musical instrument samples}. Retrieved August 2026, from \url{https://theremin.music.uiowa.edu/MIS.html}

\end{thebibliography}
